\chapter[30]{Type I Singularities and Ancient Solutions}
\label{ch:chow-iv-type-i-singularities-ancient-solutions}
\dcref{ch30}
\source{chow-et-al-ricci-flow-part-iv:page-0154}

\section{Reduced distance of Type A solutions}

\begin{definition}[Type A singular solution]
\label{def:chow-iv-type-a-singular-solution}
\source{chow-et-al-ricci-flow-part-iv:page-0154,chow-et-al-ricci-flow-part-iv:page-0155}
\dcref{ch30:30.1}
\group{Singular-Time Reduced Distance}

A finite-time singular Ricci flow
\((\mathcal M^n,g(t))\), \(t\in(\alpha,\omega)\), is \emph{Type A} if there
are \(M<\infty\) and \(r\in[1,3/2)\) such that
\[
 |\Rm|(x,t)\leq \frac{M}{(\omega-t)^r}
 \qquad\text{on }\mathcal M\times(\alpha,\omega).
\]
The case \(r=1\) is the Type I condition.
\end{definition}

\begin{lemma}[Shi's estimates for Type A solutions]
\label{lem:chow-iv-shi-estimates-type-a}
\source{chow-et-al-ricci-flow-part-iv:page-0155}
\uses{def:chow-iv-type-a-singular-solution}
\dcref{ch30:30.2}
\group{Singular-Time Reduced Distance}


If a complete Type A solution has constant \(M\) and exponent \(r\), then for
all \(k,\ell\geq0\) and
\(\eta\in(0,(\omega-\alpha)/2]\) there is
\(C_{n,k,\ell,\eta,M}<\infty\) such that
\[
 \left|\frac{\partial^k}{\partial t^k}\nabla^\ell\Rm\right|(x,t)
 \leq \frac{C_{n,k,\ell,\eta,M}}
 {(\omega-t)^{(1+k+\ell/2)r}}
\]
on \(\mathcal M\times[\alpha+\eta,\omega)\).  In particular,
\(|\nabla R|\leq C(\omega-t)^{-3r/2}\) there.
\end{lemma}

\begin{definition}[Reduced distance based at the singular time]
\label{def:chow-iv-reduced-distance-singular-time}
\source{chow-et-al-ricci-flow-part-iv:page-0156}
\dcref{ch30:30.3}
\group{Singular-Time Reduced Distance}

Let \((p_i,t_i)\) satisfy \(t_i\nearrow\omega\).  A subsequential limit on
\(\mathcal M\times(\alpha,\omega)\) of the ordinary reduced distances
\(\ell_{p_i,t_i}\) is called a \emph{reduced distance based at the singular
time}.
\end{definition}

For a pointed Cheeger--Gromov convergent sequence
\((\mathcal M_i,g_i(t),p_i)\to(\mathcal M_\infty,g_\infty(t),p_\infty)\),
write \(\Phi_i\) for the convergence embeddings and
\[
 \widetilde L_i(q,t)=L^{g_i}_{p_i,\omega_i}(\Phi_i(q),t),
 \qquad
 \widetilde\ell_i(q,t)=\ell^{g_i}_{p_i,\omega_i}(\Phi_i(q),t).
\]

\begin{lemma}[Uniform upper bound for \(\widetilde L_i\) on compact sets]
\label{lem:chow-iv-type-a-l-distance-upper-bound}
\source{chow-et-al-ricci-flow-part-iv:page-0157,chow-et-al-ricci-flow-part-iv:page-0158,chow-et-al-ricci-flow-part-iv:page-0159}
\uses{def:chow-iv-type-a-singular-solution}
\dcref{ch30:30.4}
\group{Singular-Time Reduced Distance}


Assume uniformly that
\(|\Rm_{g_i}|\leq M(\omega_i-t)^{-r}\), where \(r\in[1,3/2)\), and fix
\(t_0\in(\alpha_\infty,\omega_\infty)\).  If \(r>1\), then for sufficiently
large \(i\),
\[
\begin{split}
 \widetilde L_i(q,t)\leq{}&
 2\exp\!\left(\frac{2nM}{r-1}
 \left(\frac{\omega_\infty-t}{2}\right)^{-(r-1)}\right)
 \sqrt{\omega_\infty-\alpha_\infty}
 \frac{d_{g_\infty(t_0)}(q,p_\infty)^2}{\omega_\infty-t}\\
 &+\frac{n^2M}{3/2-r}(\omega_\infty-t)^{3/2-r}.
\end{split}
\]
If \(r=1\), the exponential factor is replaced by
\(2(2(\omega_\infty-t_0)/(\omega_\infty-t))^{2nM}\), and the final term
may be taken to be \(2n^2M(\omega_\infty-t)^{1/2}\).
\end{lemma}

\begin{corollary}[Compact-set bound for \(\widetilde L_i\)]
\label{cor:chow-iv-type-a-l-distance-compact-bound}
\source{chow-et-al-ricci-flow-part-iv:page-0157}
\uses{lem:chow-iv-type-a-l-distance-upper-bound}
\dcref{ch30:30.5}
\group{Singular-Time Reduced Distance}


Under the hypotheses of Lemma~\ref{lem:chow-iv-type-a-l-distance-upper-bound},
for every compact \(K\subset\mathcal M_\infty\), every
\(\varepsilon>0\), and every
\(t_0<\omega_\infty-\varepsilon\), the functions \(\widetilde L_i\) are
uniformly bounded on \(K\times[t_0,\omega_\infty-\varepsilon]\) for all
sufficiently large \(i\).
\end{corollary}

\begin{lemma}[Uniform gradient bounds for \(\widetilde L_i\)]
\label{lem:chow-iv-type-a-l-distance-gradient-bound}
\source{chow-et-al-ricci-flow-part-iv:page-0159,chow-et-al-ricci-flow-part-iv:page-0160,chow-et-al-ricci-flow-part-iv:page-0161}
\uses{lem:chow-iv-shi-estimates-type-a, lem:chow-iv-type-a-l-distance-upper-bound}
\dcref{ch30:30.7}
\group{Singular-Time Reduced Distance}


Under the hypotheses of Lemma~\ref{lem:chow-iv-type-a-l-distance-upper-bound},
for every compact \(K\subset\mathcal M_\infty\) and
\(\varepsilon\in(0,(\omega_\infty-\alpha_\infty)/2)\), there is a constant
independent of \(i\) such that
\[
 |\nabla\widetilde L_i|_{\Phi_i^*g_i}(q,t)\leq C
\]
on \(K\times[\alpha_\infty+\varepsilon,
\omega_\infty-\varepsilon]\) for all sufficiently large \(i\).
\end{lemma}

\begin{lemma}[Uniform time-derivative bounds for \(\widetilde L_i\)]
\label{lem:chow-iv-type-a-l-distance-time-bound}
\source{chow-et-al-ricci-flow-part-iv:page-0161,chow-et-al-ricci-flow-part-iv:page-0162}
\uses{lem:chow-iv-type-a-l-distance-gradient-bound}
\dcref{ch30:30.8}
\group{Singular-Time Reduced Distance}


On the same compact space-time sets as in
Lemma~\ref{lem:chow-iv-type-a-l-distance-gradient-bound},
\[
 \left|\frac{\partial\widetilde L_i}{\partial t}\right|\leq C
\]
uniformly for all sufficiently large \(i\).
\end{lemma}

\begin{theorem}[Reduced distance based at the singular time for a sequence of solutions]
\label{thm:chow-iv-singular-time-reduced-distance-limit}
\source{chow-et-al-ricci-flow-part-iv:page-0162,chow-et-al-ricci-flow-part-iv:page-0163,chow-et-al-ricci-flow-part-iv:page-0164,chow-et-al-ricci-flow-part-iv:page-0165}
\uses{def:chow-iv-reduced-distance-singular-time, cor:chow-iv-type-a-l-distance-compact-bound, lem:chow-iv-type-a-l-distance-gradient-bound, lem:chow-iv-type-a-l-distance-time-bound}
\dcref{ch30:30.10}
\group{Singular-Time Reduced Distance}


Suppose complete bounded-curvature pointed Ricci flows
\((\mathcal M_i,g_i(t),p_i)\), \(t\in[\alpha_i,\omega_i]\), converge smoothly
in pointed Cheeger--Gromov sense to
\((\mathcal M_\infty,g_\infty(t),p_\infty)\), with
\(\omega_i\nearrow\omega_\infty\), and satisfy the uniform Type A estimate.
Then a subsequence of \(\widetilde\ell_i\) converges in
\(C^\beta_{\mathrm{loc}}\), for every \(\beta\in(0,1)\), to a locally
Lipschitz reduced distance \(\ell_\infty\).  In both the support and weak
senses,
\[
 -\partial_t\ell_\infty-\Delta_{g_\infty}\ell_\infty
 +|\nabla\ell_\infty|^2-R_{g_\infty}
 +\frac{n}{2(\omega_\infty-t)}\geq0.
\]
\end{theorem}

\begin{corollary}[Reduced distance based at the singular time for a Type A solution]
\label{cor:chow-iv-type-a-singular-time-reduced-distance}
\source{chow-et-al-ricci-flow-part-iv:page-0162,chow-et-al-ricci-flow-part-iv:page-0163}
\uses{thm:chow-iv-singular-time-reduced-distance-limit}
\dcref{ch30:30.11}
\group{Singular-Time Reduced Distance}


For a complete Type A singular solution on \([0,T)\), every \(p\in\mathcal M\)
and sequence \(T_i\nearrow T\) admit a subsequence for which
\(\ell_{p,T_i}\) converges locally in \(C^\beta\), \(0<\beta<1\), to a
locally Lipschitz reduced distance \(\ell_{p,T}\).  It satisfies
\[
 -\partial_t\ell_{p,T}-\Delta\ell_{p,T}+|\nabla\ell_{p,T}|^2-R
 +\frac{n}{2(T-t)}\geq0
\]
in the support and weak senses.
\end{corollary}

\begin{remark}[Strong convergence of reduced-distance gradients]
\label{rem:chow-iv-strong-reduced-distance-gradient-convergence}
\source{chow-et-al-ricci-flow-part-iv:page-0165}
\uses{thm:chow-iv-singular-time-reduced-distance-limit}
\dcref{ch30:30.12}
\group{Singular-Time Reduced Distance}


The proof gives
\(|\nabla\widetilde\ell_i|^2_{\Phi_i^*g_i}
\to|\nabla\ell_\infty|^2_{g_\infty}\) in distributions.  Together with
weak \(W^{1,2}_{\mathrm{loc}}\) convergence, this identifies the limiting
quadratic gradient term in the reduced-distance inequality.
\end{remark}

\section{Reduced volume at the singular time}

For a singular-time reduced distance, define
\[
 \widetilde V_\infty(t)=
 \int_{\mathcal M_\infty}(4\pi(\omega_\infty-t))^{-n/2}
 e^{-\ell_\infty(x,t)}\,d\mu_{g_\infty(t)}(x).
\]

\begin{theorem}[Reduced-volume monotonicity based at the singular time]
\label{thm:chow-iv-singular-time-reduced-volume-monotonicity}
\source{chow-et-al-ricci-flow-part-iv:page-0166,chow-et-al-ricci-flow-part-iv:page-0172,chow-et-al-ricci-flow-part-iv:page-0173,chow-et-al-ricci-flow-part-iv:page-0174,chow-et-al-ricci-flow-part-iv:page-0175}
\uses{thm:chow-iv-singular-time-reduced-distance-limit, prop:chow-iv-limit-reduced-distance-estimates}
\dcref{ch30:30.13}
\group{Singular-Time Reduced Volume}


Under the hypotheses of
Theorem~\ref{thm:chow-iv-singular-time-reduced-distance-limit} with \(r=1\),
the reduced volume is well defined, satisfies \(\widetilde V_\infty\leq1\),
and is differentiable with \(d\widetilde V_\infty/dt\geq0\).  If it has
equal values at two times \(t_1<t_2\), then on that time interval
\[
 \Ric_{g_\infty}+\Hess_{g_\infty}\ell_\infty
 -\frac{1}{2(\omega_\infty-t)}g_\infty=0;
\]
thus the limit is a shrinking gradient Ricci soliton with potential
\(\ell_\infty\).
\end{theorem}

\begin{corollary}[Reduced-volume monotonicity for a Type I solution]
\label{cor:chow-iv-type-i-singular-time-reduced-volume}
\source{chow-et-al-ricci-flow-part-iv:page-0166}
\uses{cor:chow-iv-type-a-singular-time-reduced-distance, thm:chow-iv-singular-time-reduced-volume-monotonicity}
\dcref{ch30:30.14}
\group{Singular-Time Reduced Volume}


For a complete Type I singular solution on \([0,T)\), the reduced volume
associated with \(\ell_{p,T}\) is at most one and nondecreasing.  Equality at
two times makes
\((\mathcal M,g(t),\ell_{p,T}(t),-1/(T-t))\) a shrinking gradient Ricci
soliton on the intervening interval.
\end{corollary}

\begin{lemma}[The reduced distance of a Type I solution has at most quadratic growth in space]
\label{lem:chow-iv-type-i-reduced-distance-quadratic-upper-growth}
\source{chow-et-al-ricci-flow-part-iv:page-0167,chow-et-al-ricci-flow-part-iv:page-0168,chow-et-al-ricci-flow-part-iv:page-0169}
\dcref{ch30:30.15}
\group{Singular-Time Reduced Volume}

Let a complete bounded-curvature Ricci flow on \([\alpha,\omega]\) satisfy
\(|\Rm|\leq M/(\omega-t)\), and put \(\ell=\ell_{p,\omega}\).  For every
\(\varepsilon\in(0,(\omega-\alpha)/2)\) there are \(A,B<\infty\), depending
only on \(n,\varepsilon,M\), such that
\[
 \ell(x,t)\leq
 \frac{B d_{g(t)}(q,x)^2}{2(\omega-t)}+\frac AB+2\ell(q,t)
\]
for \(t\geq\alpha+\varepsilon\).  In particular,
\[
 \ell(x,t)\leq
 \frac{B d_{g(t)}(p,x)^2}{2(\omega-t)}+\frac AB+2n^2M.
\]
\end{lemma}

\begin{corollary}[Gradient and time-derivative bounds for Type I reduced distance]
\label{cor:chow-iv-type-i-reduced-distance-derivative-bounds}
\source{chow-et-al-ricci-flow-part-iv:page-0169,chow-et-al-ricci-flow-part-iv:page-0170}
\uses{lem:chow-iv-type-i-reduced-distance-quadratic-upper-growth}
\dcref{ch30:30.16}
\group{Singular-Time Reduced Volume}


Under the hypotheses of
Lemma~\ref{lem:chow-iv-type-i-reduced-distance-quadratic-upper-growth}, there
are \(C,D,E,F<\infty\), depending only on \(n,\varepsilon,M\), such that
\[
 |\nabla\ell|(x,t)\leq \frac C{\sqrt{\omega-t}}
 +\frac{D d_{g(t)}(p,x)}{\omega-t},
\qquad
 |\partial_t\ell|(x,t)\leq\frac E{\omega-t}
 +\frac{F d_{g(t)}(p,x)^2}{(\omega-t)^2}.
\]
\end{corollary}

\begin{lemma}[Lower bound for the reduced distance of a Type I solution]
\label{lem:chow-iv-type-i-reduced-distance-lower-bound}
\source{chow-et-al-ricci-flow-part-iv:page-0170,chow-et-al-ricci-flow-part-iv:page-0171,chow-et-al-ricci-flow-part-iv:page-0172}
\uses{lem:chow-iv-type-i-reduced-distance-quadratic-upper-growth, cor:chow-iv-type-i-reduced-distance-derivative-bounds}
\dcref{ch30:30.17}
\group{Singular-Time Reduced Volume}


Under the preceding Type I hypotheses, there is \(C<\infty\) such that
\[
 \sqrt{\ell(x_1,t)+A/B}+\sqrt{\ell(x_2,t)+A/B}
 \geq \frac{d_{g(t)}(x_1,x_2)}{4\sqrt B\sqrt{\omega-t}}-C.
\]
In particular,
\[
 \sqrt{\ell(x,t)+A/B}
 \geq\frac{d_{g(t)}(x,p)}{4\sqrt B\sqrt{\omega-t}}-C.
\]
\end{lemma}

\begin{proposition}[Estimates for a limiting reduced distance]
\label{prop:chow-iv-limit-reduced-distance-estimates}
\source{chow-et-al-ricci-flow-part-iv:page-0172}
\uses{lem:chow-iv-type-i-reduced-distance-quadratic-upper-growth, cor:chow-iv-type-i-reduced-distance-derivative-bounds, lem:chow-iv-type-i-reduced-distance-lower-bound, thm:chow-iv-singular-time-reduced-distance-limit}
\dcref{ch30:30.18}
\group{Singular-Time Reduced Volume}


Under the Type I hypotheses of
Theorem~\ref{thm:chow-iv-singular-time-reduced-distance-limit}, for every
positive time margin there is \(C<\infty\) such that
\[
 \frac{d_{g_\infty(t)}(x,p_\infty)^2}{C(\omega_\infty-t)}-C
 \leq\ell_\infty(x,t)\leq
 \frac{C d_{g_\infty(t)}(x,p_\infty)^2}{\omega_\infty-t}+C,
\]
\[
 |\nabla\ell_\infty|\leq C\left((\omega_\infty-t)^{-1/2}
 +\frac{d_{g_\infty(t)}(x,p_\infty)}{\omega_\infty-t}\right),
\]
and \(|\partial_t\ell_\infty|\) has the corresponding quadratic growth.
The same estimates hold for a single Type I solution.
\end{proposition}

\section{Shrinker singularity models}

\begin{proposition}[Limits of Type I solutions with fixed basepoints are possibly flat shrinkers]
\label{prop:chow-iv-fixed-basepoint-type-i-limits-shrinkers}
\source{chow-et-al-ricci-flow-part-iv:page-0176,chow-et-al-ricci-flow-part-iv:page-0177}
\uses{cor:chow-iv-type-i-singular-time-reduced-volume, prop:chow-iv-limit-reduced-distance-estimates}
\dcref{ch30:30.19}
\group{Type I Singularity Models}


Let \((\mathcal M,g(t))\), \(t\in[0,\omega)\), be a Type I singular
solution on a closed manifold.  For any fixed \(p\in\mathcal M\) and any
\(\lambda_i\to\infty\), every pointed limit of
\[
 g_i(t)=\lambda_i g(\omega+\lambda_i^{-1}t),\qquad t<0,
\]
based at \(p\), is a complete bounded-curvature shrinking gradient Ricci
soliton.  The limit may be flat.
\end{proposition}

\begin{remark}[Flat fixed-basepoint limits]
\label{rem:chow-iv-flat-fixed-basepoint-type-i-limit}
\source{chow-et-al-ricci-flow-part-iv:page-0177}
\uses{prop:chow-iv-fixed-basepoint-type-i-limits-shrinkers}
\dcref{ch30:30.20}
\group{Type I Singularity Models}


A fixed basepoint can miss the concentrating curvature.  In dimension three,
if its rescaled curvature tends to zero, the limiting shrinker is flat by the
strong maximum principle and uniqueness of complete bounded-curvature Ricci
flows.
\end{remark}

\begin{definition}[Set of Type I singular points]
\label{def:chow-iv-type-i-singular-point-set}
\source{chow-et-al-ricci-flow-part-iv:page-0178}
\dcref{ch30:30.21}
\group{Type I Singularity Models}

For a finite-time singular solution define
\[
 \Sigma_I=\left\{p^*:\ \exists\,t_i\nearrow\omega, p_i\to p^*,
 \quad\liminf_i(\omega-t_i)|\Rm|(p_i,t_i)>0\right\}.
\]
Thus \(\Sigma_I\) consists of limits of Type I essential points.
\end{definition}

\begin{theorem}[Limits based in \(\Sigma_I\) are nonflat shrinkers]
\label{thm:chow-iv-type-i-singular-point-limits-nonflat}
\source{chow-et-al-ricci-flow-part-iv:page-0178}
\uses{def:chow-iv-type-i-singular-point-set, prop:chow-iv-fixed-basepoint-type-i-limits-shrinkers}
\dcref{ch30:30.22}
\group{Type I Singularity Models}


For a Type I singular solution on a closed manifold, every fixed-basepoint
limit based at \(p^*\in\Sigma_I\), for any dilation sequence tending to
infinity, is a non-Ricci-flat shrinking gradient Ricci soliton.
\end{theorem}

\begin{theorem}[Type I singular solutions have unbounded scalar curvature]
\label{thm:chow-iv-type-i-singular-scalar-curvature-unbounded}
\source{chow-et-al-ricci-flow-part-iv:page-0179}
\uses{thm:chow-iv-type-i-singular-point-limits-nonflat}
\dcref{ch30:30.23}
\group{Type I Singularity Models}


Every finite-time Type I singular solution on a closed manifold satisfies
\[
 \sup_{\mathcal M\times[0,T)}R=\infty.
\]
\end{theorem}

\begin{definition}[Sets of singular points]
\label{def:chow-iv-singular-point-sets}
\source{chow-et-al-ricci-flow-part-iv:page-0179}
\dcref{ch30:30.24}
\group{Type I Singularity Models}

For a finite-time singular solution set
\[
 \Sigma_R=\{p:\liminf_{t\nearrow\omega}(\omega-t)R(p,t)>0\}
\]
and
\[
 \Sigma=\{p:\exists,t_i\nearrow\omega, p_i\to p,
 |\Rm|(p_i,t_i)\to\infty\}.
\]
One always has \(\Sigma_R\subseteq\Sigma_I\subseteq\Sigma\).
\end{definition}

\begin{theorem}[Equality of the sets of singular points]
\label{thm:chow-iv-type-i-singular-point-sets-equal}
\source{chow-et-al-ricci-flow-part-iv:page-0179,chow-et-al-ricci-flow-part-iv:page-0180}
\uses{def:chow-iv-type-i-singular-point-set, def:chow-iv-singular-point-sets, prop:chow-iv-fixed-basepoint-type-i-limits-shrinkers, thm:chow-iv-type-i-singular-point-limits-nonflat}
\dcref{ch30:30.25}
\group{Type I Singularity Models}


For a Type I singular solution on a closed manifold,
\[
 \Sigma_R=\Sigma_I=\Sigma.
\]
\end{theorem}

\begin{corollary}[Distinguishing singular and nonsingular points]
\label{cor:chow-iv-distinguish-type-i-singular-points}
\source{chow-et-al-ricci-flow-part-iv:page-0180}
\uses{thm:chow-iv-type-i-singular-point-sets-equal}
\dcref{ch30:30.26}
\group{Type I Singularity Models}


At every point \(p\) of a closed Type I singular solution, either
\(\liminf_{t\nearrow\omega}(\omega-t)R(p,t)>0\), or curvature remains
uniformly bounded along every sequence \((p_i,t_i)\to(p,\omega)\).
\end{corollary}

\begin{theorem}[Limits based in \(\Sigma\) are nonflat shrinkers]
\label{thm:chow-iv-singular-point-curvature-scale-limits-nonflat}
\source{chow-et-al-ricci-flow-part-iv:page-0180}
\uses{thm:chow-iv-type-i-singular-point-sets-equal, thm:chow-iv-type-i-singular-point-limits-nonflat}
\dcref{ch30:30.27}
\group{Type I Singularity Models}


Let \(p\in\Sigma\) for a Type I singular solution on a closed manifold.  For
any \(t_i\nearrow\omega\), the curvature-scale rescalings
\[
 \widehat g_i(t)=|\Rm|(p,t_i)
 g\bigl(t_i+|\Rm|(p,t_i)^{-1}t\bigr)
\]
subconverge to a non-Ricci-flat shrinking gradient Ricci soliton.
\end{theorem}

\section{Type I ancient solutions}

For an ancient solution and positive sequences
\(\tau_i^+\searrow0\), \(\tau_i^-\nearrow\infty\), write
\[
 g_i^\pm(t)=(\tau_i^\pm)^{-1}g(\tau_i^\pm t),\qquad t<0.
\]

\begin{lemma}[Characterization of canonical form for shrinkers]
\label{lem:chow-iv-shrinker-canonical-form-characterization}
\source{chow-et-al-ricci-flow-part-iv:page-0181,chow-et-al-ricci-flow-part-iv:page-0182}
\dcref{ch30:30.29}
\group{Type I Ancient Solutions}

Let \((\mathcal M,g(t),f(t),1/t)\), \(t<0\), be a complete
bounded-curvature shrinking gradient Ricci soliton.  It is in canonical form,
\[
 g(t)=-t\,\varphi(t)^*g(-1),\qquad
 f(t)=f(-1)\circ\varphi(t),
\]
if and only if \(\partial_t f=|\nabla f|^2\).
\end{lemma}

\begin{theorem}[Backward and forward asymptotic shrinkers for Type I ancient solutions]
\label{thm:chow-iv-type-i-ancient-asymptotic-shrinkers}
\source{chow-et-al-ricci-flow-part-iv:page-0182,chow-et-al-ricci-flow-part-iv:page-0183}
\uses{cor:chow-iv-type-a-singular-time-reduced-distance, prop:chow-iv-limit-reduced-distance-estimates, lem:chow-iv-shrinker-canonical-form-characterization}
\dcref{ch30:30.31}
\group{Type I Ancient Solutions}


Let a \(\kappa\)-noncollapsed ancient Ricci flow satisfy
\(|\Rm|(x,t)\leq C/|t|\).  For every basepoint and every pair of sequences
\(\tau_i^+\searrow0\), \(\tau_i^-\nearrow\infty\), subsequences of the
rescaled flows converge to complete shrinking gradient Ricci solitons.  Both
limits are \(\kappa\)-noncollapsed and in canonical form, but either may be
Euclidean.
\end{theorem}

\begin{proposition}[Backward limits with fixed basepoint of a shrinker]
\label{prop:chow-iv-shrinker-backward-limit-isometric}
\source{chow-et-al-ricci-flow-part-iv:page-0184}
\uses{lem:chow-iv-shrinker-canonical-form-characterization}
\dcref{ch30:30.32}
\group{Type I Ancient Solutions}


Let a complete bounded-curvature shrinking gradient Ricci soliton be in
canonical form.  For any fixed basepoint and \(\tau_i\to\infty\), every
Cheeger--Gromov convergent subsequence of its backward dilations has limit
isometric to the original soliton, based at a critical point of its potential.
\end{proposition}

\begin{theorem}[Rigidity when forward and backward asymptotic solitons agree]
\label{thm:chow-iv-equal-asymptotic-shrinkers-rigidity}
\source{chow-et-al-ricci-flow-part-iv:page-0184,chow-et-al-ricci-flow-part-iv:page-0185,chow-et-al-ricci-flow-part-iv:page-0186}
\uses{thm:chow-iv-type-i-ancient-asymptotic-shrinkers, prop:chow-iv-shrinker-backward-limit-isometric, cor:chow-iv-type-i-singular-time-reduced-volume}
\dcref{ch30:30.33}
\group{Type I Ancient Solutions}


If the forward and backward asymptotic solitons in
Theorem~\ref{thm:chow-iv-type-i-ancient-asymptotic-shrinkers} are isometric
as Ricci flows, then each is isometric to the original ancient solution.
\end{theorem}

\begin{theorem}[Existence of asymptotic shrinkers for Type I ancient solutions]
\label{thm:chow-iv-type-i-ancient-nonflat-backward-shrinker}
\source{chow-et-al-ricci-flow-part-iv:page-0186}
\dcref{ch30:30.34}
\group{Type I Ancient Solutions}

For every complete \(\kappa\)-noncollapsed Type I ancient Ricci flow, some
backward limit is a nonflat shrinker.
\end{theorem}

\begin{theorem}[Three-dimensional Type II ancient solutions have asymptotic shrinkers]
\label{thm:chow-iv-three-dimensional-type-ii-asymptotic-shrinker}
\source{chow-et-al-ricci-flow-part-iv:page-0186,chow-et-al-ricci-flow-part-iv:page-0187}
\dcref{ch30:30.36}
\group{Type I Ancient Solutions}

Every three-dimensional \(\kappa\)-noncollapsed Type II ancient Ricci flow
with bounded curvature has a backward asymptotic shrinker.
\end{theorem}

\section{Nongradient shrinkers}

\begin{lemma}[Bounds for vector fields of shrinking Ricci solitons]
\label{lem:chow-iv-shrinker-vector-field-radial-bound}
\source{chow-et-al-ricci-flow-part-iv:page-0187}
\dcref{ch30:30.37}
\group{Nongradient Shrinkers}

Let \((\mathcal M^n,g,X,-1)\) be a complete noncompact shrinking Ricci
soliton structure with bounded Ricci curvature, so that
\[
 2R_{ij}+\nabla_iX_j+\nabla_jX_i-g_{ij}=0.
\]
For \(r(x)=d_g(x,O)\), there is \(C<\infty\) such that
\[
 \langle X,\nabla r\rangle(x)\geq\frac{r(x)}2-C
\]
away from \(O\).
\end{lemma}

\begin{lemma}[Completeness of nongradient shrinking-soliton vector fields]
\label{lem:chow-iv-shrinker-vector-field-complete}
\source{chow-et-al-ricci-flow-part-iv:page-0187,chow-et-al-ricci-flow-part-iv:page-0188}
\dcref{ch30:30.38}
\group{Nongradient Shrinkers}

If \((\mathcal M^n,g,X,-1)\) is a complete noncompact shrinking Ricci
soliton structure with bounded Ricci curvature, then \(X\) is a complete
vector field.
\end{lemma}

\begin{theorem}[Bounded-curvature shrinkers are gradient and noncollapsed]
\label{thm:chow-iv-bounded-curvature-shrinker-gradient-noncollapsed}
\source{chow-et-al-ricci-flow-part-iv:page-0189}
\uses{lem:chow-iv-shrinker-vector-field-radial-bound, lem:chow-iv-shrinker-vector-field-complete, thm:chow-iv-type-i-ancient-asymptotic-shrinkers, prop:chow-iv-shrinker-backward-limit-isometric}
\dcref{ch30:30.40}
\group{Nongradient Shrinkers}


Let \((\mathcal M,g(t),X(t),1/t)\), \(t<0\), be a complete noncompact
bounded-curvature shrinking Ricci soliton in canonical form.  Then there is a
normalized potential \(f\) making it a complete shrinking gradient Ricci
soliton.  Moreover it is \(\kappa\)-noncollapsed below all scales for a
\(\kappa>0\) depending only on \(n\) and \(\int e^{-f}\,d\mu\).
\end{theorem}

\begin{conjecture}[Problem 30.28: Type A shrinker models]
\label{conj:chow-iv-type-a-nonflat-shrinker-model}
\dcref{ch30:30.28}
\group{Open Problems}
\source{chow-et-al-ricci-flow-part-iv:page-0180}
Does every Type A singular solution on a closed manifold admit a nonflat
shrinking gradient Ricci soliton as a singularity model?
\end{conjecture}

\begin{conjecture}[Problem 30.35: Type II asymptotic shrinkers]
\label{conj:chow-iv-type-ii-ancient-backward-shrinker}
\dcref{ch30:30.35}
\group{Open Problems}
\source{chow-et-al-ricci-flow-part-iv:page-0186}
Does every \(\kappa\)-noncollapsed Type II ancient solution have a backward
asymptotic shrinker?
\end{conjecture}

\begin{example}[Exercise 30.9: explicit first-derivative estimates]
\label{ex:chow-iv-explicit-l-distance-derivative-bounds}
\dcref{ch30:30.9}
\group{Exercises}
\source{chow-et-al-ricci-flow-part-iv:page-0162}
\uses{lem:chow-iv-type-a-l-distance-gradient-bound,lem:chow-iv-type-a-l-distance-time-bound}
Derive explicit bounds for \(|\nabla\widetilde L_i|\) and
\(|\partial_t\widetilde L_i|\) in terms of time and distance from a fixed
point, analogous to Lemma~\ref{lem:chow-iv-type-a-l-distance-upper-bound}.
\end{example}

\begin{example}[Exercise 30.6: optimizing the test path]
\label{ex:chow-iv-optimize-type-a-test-path}
\dcref{ch30:30.6}
\group{Exercises}
\source{chow-et-al-ricci-flow-part-iv:page-0159}
\uses{lem:chow-iv-type-a-l-distance-upper-bound}
Determine whether the estimate for \(r>1\) can be qualitatively strengthened
by keeping the test path at \(p_i\) only on
\([\omega_\infty-\delta,\omega_i]\), with
\(\delta\in(0,\omega_\infty-t)\) chosen optimally.
\end{example}

\begin{example}[Exercise 30.30: normalization in canonical form]
\label{ex:chow-iv-shrinker-canonical-normalization}
\dcref{ch30:30.30}
\group{Exercises}
\source{chow-et-al-ricci-flow-part-iv:page-0182}
\uses{lem:chow-iv-shrinker-canonical-form-characterization}
If \(f\) is normalized by \(R+|\nabla f|^2-f=0\) at time \(-1\), then
show that
\[
 R+|\nabla f|^2+\frac ft=0
\]
for every \(t<0\).  Indeed, the time derivative of the left-hand side is
\(-1/t\) times that expression.
\end{example}

\begin{example}[Exercise 30.39: evolution of the canonical soliton vector field]
\label{ex:chow-iv-shrinker-vector-field-evolution}
\dcref{ch30:30.39}
\group{Exercises}
\source{chow-et-al-ricci-flow-part-iv:page-0188}
\uses{lem:chow-iv-shrinker-vector-field-complete}
For the canonical family
\[
 g(t)=-t\,\varphi(t)^*g,
 \qquad X(t)=-t^{-1}\varphi(t)^*X,
\]
compute \(\partial_tX(t)\).
\end{example}
